% ---------------------------------------------------------------------
\chapter{Befehle lesen und verstehen}
\label{kap:lesen}

\begin{lernziele}
\begin{itemize}[nosep]
  \item wissen, was ein Terminal ist und was der aktuelle Ordner bedeutet,
  \item jeden Befehl in seine Bestandteile zerlegen können,
  \item den Unterschied zwischen \cmd{-x} und \cmd{--wort} begründen können,
  \item sich zu jedem Befehl selbst Hilfe holen können.
\end{itemize}
\end{lernziele}

\section{Das Terminal}

Git wird meist über ein \begriff{Terminal} bedient, ein Fenster, in das man Befehle als Text
eingibt. Unter Windows sind das zum Beispiel \emph{PowerShell} oder \emph{Git Bash}, in VS Code
das Terminal am unteren Rand. Grafische Oberflächen für Git gibt es auch, sie bauen aber auf
denselben Befehlen auf. Wer die Befehle versteht, versteht auch jede Oberfläche.

Das Terminal zeigt vor dem Cursor eine Zeile, die \begriff{Eingabeaufforderung}
(\emph{prompt}). Sie verrät, wo man sich befindet:

\begin{lstlisting}[style=ausgabe]
PS D:\github\hslu-bachelor-ai-ml\2026-FA\ADS>
bahram@PC MINGW64 /d/github/hslu-bachelor-ai-ml (main)
\end{lstlisting}

Die erste Zeile stammt aus PowerShell, die zweite aus Git Bash. Git Bash zeigt in Klammern
zusätzlich den aktuellen Branch an, hier \cmd{main}.

\subsection*{Der aktuelle Ordner}

Jedes Terminal „steht“ in einem Ordner, dem \begriff{aktuellen Arbeitsverzeichnis}. Alle
Pfade, die man angibt, gelten \textbf{relativ zu diesem Ordner}. Das ist der häufigste Grund,
warum ein Befehl „nicht funktioniert“: Man steht im falschen Ordner.

\begin{lstlisting}[style=befehle]
cd D:\github\hslu-bachelor-ai-ml       # in den Ordner wechseln (change directory)
cd 2026-FA\ADS                         # relativ: einen Unterordner tiefer
cd ..                                  # einen Ordner nach oben
pwd                                    # zeigen, wo man steht (print working dir.)
\end{lstlisting}

Git-Befehle wirken immer auf das Repository, in dem der aktuelle Ordner liegt. Git sucht
dazu vom aktuellen Ordner aus nach oben, bis es den Ordner \cmd{.git} findet. Deshalb
funktionieren Git-Befehle in jedem Unterordner eines Repositorys.

\section{Aufbau eines Befehls}

Alle Git-Befehle folgen demselben Muster:

\begin{center}
\begin{tikzpicture}[
    teil/.style={font=\ttfamily\large, inner sep=2pt},
    erkl/.style={font=\small, align=center},
    klammer/.style={decorate, decoration={brace, amplitude=4pt, mirror}, thick, hsluBlau}
  ]
  \node[teil] (p)  at (0,0)    {git};
  \node[teil] (u)  at (2.2,0)  {diff};
  \node[teil] (o1) at (4.6,0)  {--cached};
  \node[teil] (o2) at (7.0,0)  {--stat};
  \node[teil] (t)  at (9.0,0)  {--};
  \node[teil] (a)  at (11.2,0) {datei.py};
  \draw[klammer] (p.south west)  -- (p.south east)  node[midway, below=6pt, erkl] {Programm};
  \draw[klammer] (u.south west)  -- (u.south east)  node[midway, below=6pt, erkl] {Unterbefehl:\\was getan wird};
  \draw[klammer] (o1.south west) -- (o2.south east) node[midway, below=6pt, erkl] {Optionen:\\wie es getan wird};
  \draw[klammer] (t.south west)  -- (t.south east)  node[midway, below=6pt, erkl] {Trenner};
  \draw[klammer] (a.south west)  -- (a.south east)  node[midway, below=6pt, erkl] {Argument:\\worauf};
\end{tikzpicture}
\end{center}

\begin{description}[style=nextline, leftmargin=1.2cm]
  \item[Programm] \cmd{git} oder \cmd{gh}. Das Terminal startet dieses Programm.
  \item[Unterbefehl] die Tätigkeit: \cmd{status}, \cmd{add}, \cmd{commit}, \ldots{} Git hat
    über hundert Unterbefehle. Für den Alltag genügen etwa fünfzehn.
  \item[Optionen] verändern \emph{wie} der Unterbefehl arbeitet. Sie beginnen mit einem
    oder zwei Strichen. Die meisten sind freiwillig.
  \item[Argumente] sagen, \emph{worauf} sich der Befehl bezieht, zum Beispiel Dateien,
    Ordner oder Branch-Namen.
\end{description}

\section{Ein Strich oder zwei Striche}

Die Schreibweise der Optionen ist keine Eigenheit von Git, sondern eine Konvention, die
fast alle Kommandozeilen-Programme seit den 1990er-Jahren teilen:

\begin{center}
\renewcommand{\arraystretch}{1.3}
\begin{tabularx}{\linewidth}{@{}l>{\ttfamily}l>{\raggedright\arraybackslash}X@{}}
  \toprule
  \textbf{Form} & \normalfont\textbf{Beispiel} & \textbf{Zweck} \\
  \midrule
  ein Strich, ein Buchstabe & -m, -s, -b & Kurzform: schnell zu tippen, für den Alltag \\
  zwei Striche, ganzes Wort & --message, --short & Langform: selbsterklärend, für Skripte und Dokumentation \\
  \bottomrule
\end{tabularx}
\end{center}

Viele Optionen gibt es in beiden Formen: \cmd{-m} und \cmd{--message} bedeuten dasselbe.

Die zwei Striche haben einen technischen Grund. Kurzformen dürfen \textbf{zusammengezogen}
werden: \cmd{-sb} bedeutet dasselbe wie \cmd{-s -b}. Würde man ganze Wörter mit nur einem
Strich schreiben, wäre \cmd{-stat} mehrdeutig: die Option \cmd{stat}, oder die vier
Optionen \cmd{-s -t -a -t}? Der doppelte Strich beseitigt diese Mehrdeutigkeit.

\subsection*{Optionen mit Wert}

Manche Optionen brauchen einen Wert. Bei der Kurzform folgt er nach einem Leerzeichen, bei
der Langform meist nach einem Gleichheitszeichen:

\begin{lstlisting}[style=befehle]
git commit -m "Stack implementiert"     # Kurzform
git log --format="%h %s"                # Langform
\end{lstlisting}

Enthält ein Wert Leerzeichen, gehört er in \textbf{Anführungszeichen}. Das Terminal trennt
Eingaben an Leerzeichen auf. Ohne Anführungszeichen würde
\cmd{git commit -m Stack implementiert} nur \cmd{Stack} als Nachricht verwenden und
\cmd{implementiert} als Dateinamen deuten.

\subsection*{Zwei Striche ohne Wort}

Ein alleinstehendes \cmd{--} bedeutet: „Ab hier folgen keine Optionen mehr, nur noch Pfade.“
Das ist nötig, wenn ein Dateiname wie eine Option oder ein Branch aussieht:

\begin{lstlisting}[style=befehle]
git restore -- main.py         # "main.py" ist sicher eine Datei, kein Branch "main"
\end{lstlisting}

Im Alltag braucht man das selten. Es taucht aber oft in Anleitungen und Fehlermeldungen auf.

\section{Hilfe zu jedem Befehl}

Niemand kennt alle Optionen auswendig. Erfahrene Entwicklerinnen und Entwickler schlagen
ständig nach. Git liefert die Dokumentation gleich mit:

\begin{lstlisting}[style=befehle]
git diff -h                    # kurze Übersicht aller Optionen im Terminal
git help diff                  # ausführliches Handbuch im Browser
\end{lstlisting}

\begin{erfahrung}
Fehlermeldungen von Git sind ungewöhnlich hilfreich. Sie nennen fast immer die Ursache und
oft schon den passenden Befehl, besonders in Zeilen, die mit \cmd{hint:} beginnen. Wer sie
in Ruhe liest, statt sie wegzuklicken, löst die meisten Probleme selbst.
\end{erfahrung}

\begin{zusammenfassung}
\begin{itemize}[nosep]
  \item Pfade gelten relativ zum aktuellen Ordner des Terminals.
  \item Aufbau: Programm, Unterbefehl, Optionen, Argumente.
  \item \cmd{-x} ist die Kurzform, \cmd{--wort} die Langform. Kurzformen lassen sich zusammenziehen.
  \item Werte mit Leerzeichen gehören in Anführungszeichen. \cmd{--} allein trennt Optionen von Pfaden.
  \item \cmd{git <befehl> -h} und \cmd{git help <befehl>} liefern Hilfe.
\end{itemize}
\end{zusammenfassung}
